\documentclass[12pt]{article}

%packages
\usepackage[utf8]{inputenc}
\usepackage[T1]{fontenc}
\usepackage{lmodern}
\usepackage{geometry}
\usepackage{xcolor}
\usepackage{tikz}
\usepackage{hyperref}

\geometry{a4paper, margin=1in}

%colours
\definecolor{accent}{HTML}{2563EB}   % vibrant blue
\definecolor{darkbg}{HTML}{1E293B}   % dark slate
\definecolor{lightgray}{HTML}{94A3B8}

%suppress page number on title
\pagenumbering{gobble}

\begin{document}

\begin{titlepage}
\begin{center}

\vspace*{2cm}

%decorative rule
{\color{accent}\rule{\linewidth}{2pt}}

\vspace{1.5cm}

%title
{\fontsize{42}{50}\selectfont\bfseries\color{darkbg} Audio Fingerprinting\\[6pt]}

\vspace{0.8cm}

%subtitle
{\Large\color{lightgray}\textit{(PS1)}}

\vspace{1.2cm}

{\color{accent}\rule{0.4\linewidth}{1pt}}

\vspace{1.5cm}

%project info
{\LARGE\color{darkbg} Working with audio files and spectrograms}

\vspace{0.6cm}

{\large\color{lightgray} Technical Notes}

\vspace{3cm}

%author / date block
\begin{minipage}{0.6\textwidth}
\centering
{\large\color{darkbg}\textbf{Facundo Franchino} \&{} \textbf{Hans Dietrich}}

\vspace{1cm}

{\color{lightgray}\today}
\end{minipage}

\vfill

%ottom rule
{\color{accent}\rule{\linewidth}{2pt}}

\end{center}
\end{titlepage}

%start page numbering for content
\pagenumbering{arabic}


\section{Objectives}
Build the simplest possible audio fingerprinting prototype from scratch.  
The goal of this session is to load an audio file, visualise its spectrogram, 
and detect prominent spectral peaks as a first approximation to a song fingerprint.



\section{Background}
A digital audio signal is a sequence of samples taken at a fixed sample rate.  
By applying the Short-Time Fourier Transform (STFT), 
the signal can be represented as a spectrogram, 
showing how frequency content evolves over time.  
Audio fingerprinting systems such as Shazam rely 
on robust local peaks in this time--frequency representation.


\section{Setup}

We used Python 3 with \texttt{librosa}, \texttt{numpy}, \texttt{scipy}, and \texttt{matplotlib}.  
An MP3 file from the \textit{Twin Peaks} soundtrack was loaded and analysed locally.

\section{Tasks}

\subsection{Load and inspect the audio file}

Load an MP3 file, inspect its sample rate, and plot the waveform in order to understand digital audio as a sequence of samples.

\subsection{Generate and interpret the spectrogram}

Compute the STFT, convert magnitudes to decibels, and display the spectrogram using a logarithmic frequency axis.

\subsection{Detect local spectral peaks}

Apply a local maximum filter and an amplitude threshold to detect salient peaks in the spectrogram, then visualise them as candidate constellation points.

\section{Notes}

The main practical issue was choosing suitable peak-detection parameters.  
A threshold that was too low produced too many peaks, especially at low frequencies.  
Increasing the neighbourhood size, raising the threshold, and suppressing the lowest frequency bins produced a much cleaner peak map.


\end{document}